\documentclass[12pt]{article}

% --- Configuración de Fuentes (Obligatorio compilar con XeLaTeX) ---
\usepackage{fontspec}
\setmainfont{Arial} % Todo el texto general en Arial
\newfontfamily\arialblack{Arial Black} % Fuente específica para los Títulos

\usepackage[spanish]{babel}
\usepackage[margin=2.5cm, headheight=20pt]{geometry} % Margen 2.5 cm
\usepackage{setspace}
\singlespacing % Interlineado sencillo forzado
\usepackage{graphicx}
\usepackage[table]{xcolor} 
\usepackage{tikz}
\usepackage{fancyhdr}
\usepackage[most]{tcolorbox}
\usepackage[hidelinks]{hyperref}
\usepackage{titlesec}
\usepackage{amssymb}
\usepackage{enumitem}
\usepackage{multicol}
\usepackage{tabularx}
\usepackage{listings} 

% --- 1. Definición de Colores ---
\definecolor{MainBlue}{RGB}{0, 114, 187}
\definecolor{LightBlue}{RGB}{235, 245, 255}
\definecolor{AccentBlue}{RGB}{0, 180, 216}
\definecolor{DarkGray}{RGB}{50, 50, 50}
\definecolor{CodeGreen}{RGB}{0, 150, 136}
\definecolor{CodeBlue}{RGB}{0, 114, 187}
\definecolor{CodeGray}{RGB}{100, 100, 100}
\definecolor{CodeBackground}{RGB}{245, 248, 250}

% --- 2. Configuración de Listings (Códigos) ---
\renewcommand{\lstlistingname}{Estructura}
\renewcommand{\lstlistlistingname}{Índice de estructuras}

\lstset{
    language=C++,
    backgroundcolor=\color{CodeBackground},
    commentstyle=\color{CodeGray}\itshape,
    keywordstyle=\color{CodeBlue}\bfseries,
    numberstyle=\tiny\color{CodeGray},
    stringstyle=\color{CodeGreen},
    basicstyle=\ttfamily\small,
    breaklines=true,
    captionpos=b,
    numbers=left,
    numbersep=10pt,
    frame=leftline,
    rulecolor=\color{MainBlue},
    framesep=5pt
}

% --- 3. Configuración de Títulos (Formato Rúbrica) ---
% Títulos (Section): Arial-Black 16 en color negro
\titleformat{\section}{\color{black}\arialblack\fontsize{16}{19.2}\selectfont}{\thesection}{1em}{}[\titlerule]

% Subtítulos (Subsection): Arial 14 Negrita en color negro
\titleformat{\subsection}{\color{black}\fontsize{14}{16.8}\bfseries\selectfont}{\thesubsection}{1em}{}

% --- Entorno para notas ---
\newtcolorbox{infobox}[1]{
  colback=LightBlue, colframe=MainBlue, fonttitle=\bfseries,
  coltitle=white, title=#1, arc=2mm, boxrule=1pt
}

% --- Comando para listas limpias ---
\newcommand{\campo}[2]{
    \noindent\hangindent=1.2cm\makebox[0.6cm][l]{\color{MainBlue}\scriptsize$\bullet$} 
    \textbf{#1} #2 \par\vspace{0.1cm}
}

\begin{document}

% --- PORTADA ---
\begin{titlepage}
    \begin{tikzpicture}[remember picture, overlay]
        \fill[MainBlue!20] (current page.north east) circle (4cm);
        \fill[AccentBlue!30] (current page.north east) circle (2cm);
        \fill[MainBlue!10] (current page.south west) -- ([xshift=8cm]current page.south west) -- ([yshift=6cm]current page.south west) -- cycle;
        \draw[line width = 1.5pt, color=MainBlue]
        ([shift={(1cm,-1cm)}]current page.north west) rectangle ([shift={(-1cm,1cm)}]current page.south east);
    \end{tikzpicture}

    \centering
    \begin{minipage}{0.2\textwidth}
        % \includegraphics[width=\linewidth]{logo_uaslp.png} 
    \end{minipage}
    \begin{minipage}{0.58\textwidth}
        \centering
        {\scshape\Large \color{MainBlue} Universidad Autónoma de \\ San Luis Potosí \par}
        \vspace{0.2cm}
        {\color{AccentBlue}\hrule height 1.5pt}\vspace{0.2cm}
        {\footnotesize \textbf{FACULTAD DE INGENIERÍA} \par}
    \end{minipage}
    \begin{minipage}{0.2\textwidth}
        \flushright
        % \includegraphics[width=0.9\linewidth]{logo_fepzm.png}
    \end{minipage}

    \vspace{4cm}

    \begin{minipage}{0.9\textwidth}
        \centering
        \textbf{\color{MainBlue} CARRERA:} \\
        Ingeniería en Sistemas Inteligentes \\ [0.5cm]
        \textbf{\color{MainBlue} ASIGNATURA:} \\
        Lenguajes de programación \\ [0.5cm]
        \textbf{\color{MainBlue} ELABORADO POR:} \\
        \begin{tabular}{c}
            369664 Antonio Jiménez José Emilio \\
            388531 Alvarez Puente Alan Antonio
        \end{tabular} \\ [0.5cm]
        \textbf{\color{MainBlue} DOCENTE:} \\
        [0.5cm] % Espacio listo para que pongas el nombre del profe
        \textbf{\color{MainBlue} CICLO ESCOLAR:} \\
        2025 - 2026
    \end{minipage}

    \vfill
    {\large \bfseries \color{MainBlue} San Luis Potosí, S.L.P. \par}
    {\color{gray} \today}
\end{titlepage}

% --- PÁGINAS INTERNAS ---
\newpage
\pagestyle{fancy}
\fancyhf{}
\fancyhead[L]{\small \color{gray} Lenguajes de Programación}
\fancyhead[R]{\small \color{gray} Diccionario}
\fancyfoot[C]{\thepage}
\renewcommand{\headrulewidth}{0.4pt}

\tableofcontents
\newpage

\section{Estructuras de Datos y Objetos}

\phantomsection
\addcontentsline{toc}{subsection}{Entidad}
\subsection*{Descripción de la Estructura Entidad:}

La estructura es el molde inicial de nuestro Diccionario en formato de base de datos de acceso aleatorio. Se establece un alias \texttt{cadena} para asegurar que cada registro tenga el mismo tamaño.

\campo{cadena nombre:}{Identificador único de cada registro de entidad (máximo 30 caracteres).}
\campo{long atr:}{Puntero tipo \texttt{long} que conecta la entidad con su lista de atributos.}
\campo{long data:}{Dirección física en el archivo donde se almacenan los registros de datos reales.}
\campo{long sig:}{Enlace a la siguiente entidad, permitiendo una estructura de lista ligada.}

\bigskip

\begin{lstlisting}[caption=Estructura de Datos de Entidad, label=code:entidad]
#ifndef ENTIDAD_H_INCLUDED
#define ENTIDAD_H_INCLUDED
#include <iostream>
#include <cstdio>

typedef char cadena[30];
struct Entidad
{
    cadena nombre;
    long atr;
    long data;
    long sig;
};
#endif
\end{lstlisting}

% --- EJEMPLO DE DIAGRAMA ENUMERADO CON PIE DE IMAGEN ---
% Para que cumpla la regla de la imagen, descomenta esto cuando tengas tu diagrama:
% \begin{figure}[htbp]
%     \centering
%     % \includegraphics[width=0.7\textwidth]{tu_diagrama_entidad.png}
%     \caption{Diagrama de la estructura Entidad y sus relaciones.}
%     \label{fig:diagrama_entidad}
% \end{figure}

\newpage

\phantomsection
\addcontentsline{toc}{subsection}{Atributos}
\subsection*{Descripción de la Estructura Atributo:}

Esta estructura es el molde para las columnas de nuestro Diccionario. Al igual que en la entidad, se usa el alias \texttt{cadena} para mantener el tamaño fijo en el archivo binario y evitar que los registros se desfasen. En esta estructura observamos las siguientes declaraciones:

\campo{cadena nombre:}{Identificador del atributo dentro de la entidad.}
\campo{unsigned int tipo:}{Código numérico que define el tipo de dato (ej. entero, flotante, cadena).}
\campo{unsigned int tamano:}{Longitud en bytes que ocupará el valor del atributo en el registro.}
\campo{char iskp:}{Bandera (flag) que indica si el atributo funciona como llave primaria (Primary Key).}
\campo{long sig:}{Puntero lógico (offset) que enlaza con el siguiente atributo de la entidad.}
\campo{cadena descripcion:}{Información adicional o etiqueta descriptiva del campo.}
\campo{char null:}{Indicador booleano que define si el campo permite valores nulos.}

\bigskip

\begin{lstlisting}[caption=Estructura de Datos de Atributos, label=code:atributos]
#ifndef ATRIBUTO_H
#define ATRIBUTO_H

typedef char cadena [30];

typedef struct{
    cadena nombre;
    unsigned int tipo;
    unsigned int tamano;
    char iskp;
    long sig;
    cadena descripcion;
    char null;
}Atributo;

#endif
\end{lstlisting}

\newpage

\phantomsection
\addcontentsline{toc}{subsection}{CDiccionario}
\subsection*{Descripción del Objeto CDiccionario:}

La clase \texttt{CDiccionario} funciona como el administrador principal del sistema. En su parte privada, contiene las variables necesarias para controlar el flujo del archivo binario y saber exactamente en qué parte del disco estamos parados mientras el usuario navega. En esta sección observamos las siguientes declaraciones:

\campo{FILE *arch:}{Es el puntero al archivo físico donde se almacenan todas las entidades, atributos y datos.}
\campo{Entidad activa:}{Representa la entidad que el usuario está manipulando o consultando en ese momento.}
\campo{long dir:}{Guarda la dirección física (offset) actual dentro del archivo para las operaciones de lectura/escritura.}
\campo{int nAtributos:}{Es un contador que lleva el registro de cuántos atributos tiene la entidad seleccionada.}
\campo{long tambloque:}{Almacena el tamaño del bloque activo para gestionar correctamente el espacio en el archivo.}
\campo{char nombreArchivo[50]:}{Es el arreglo donde se guarda el nombre del archivo que el usuario decide abrir o crear.}

\begin{lstlisting}[caption=Estructura de Datos CDiccionario, label=code:cdiccionario]
#ifndef CDICCIONARIO_H_INCLUDED
#define CDICCIONARIO_H_INCLUDED

#include <cstdio>
#include "Atributo.h"
#include "Entidad.h"

class CDiccionario
{
private:
    FILE *arch;
    Entidad activa;
    long dir;
    int nAtributos;
    long tambloque;

    char nombreArchivo[50];
public:
    //constructor
    CDiccionario();

    //Entidades
    void Menu_Entidades();
    Entidad CapturaEntidad();
    void altaEntidad();
    long buscaEntidad(char nombre[30]);
    long getcabeceraEntidades();
    Entidad leeEntidad(long dir);
    void EscribeCabEntidades(long cab);
    long escribeEntidad(Entidad ant);
    void insertaEntidad(Entidad nvo, long dir);
    void ConsultaEntidades();
    void BajaEntidad();
    long EliminaEntidad(char nombre[30]);
    void reescribeEntidad(long dir, Entidad ant);
    void modificaEntidad();

    //Atributos
    void Menu_Atributos();
    void consultarAtributo();
    long eliminaAtributo(char cad[30]);
    void modificaAtributo();
    void AltaAtributo();
    Atributo CapturaAt();
    void insertaAtributo(Atributo nvo, long dir);
    long BuscarAtributo(char *atr);
    Atributo leeAtributo(long dir);
    long EscribeAtributo(Atributo ant);
    void reescribeAtributo(long dir, Atributo atr);
    void BajaAtributo();

    //Datos
    void Menu_Datos();

    //Principal
    void menu();
    void Nuevo();
    void abrir();

};

#endif
\end{lstlisting}

\newpage
\section[Funciones]{Funciones \hfill \fontsize{14}{17}\bfseries\selectfont Entidades}

\phantomsection
\addcontentsline{toc}{subsection}{Menu Entidades}
\subsection*{Menu Entidades}

\phantomsection
\addcontentsline{toc}{subsection}{Captura Entidad}
\subsection*{Captura Entidad}

\phantomsection
\addcontentsline{toc}{subsection}{Alta Entidad}
\subsection*{Alta Entidad}

\phantomsection
\addcontentsline{toc}{subsection}{Busca Entidad}
\subsection*{Busca Entidad}

\phantomsection
\addcontentsline{toc}{subsection}{Get Cabecera Entidades}
\subsection*{Get Cabecera Entidades}

\phantomsection
\addcontentsline{toc}{subsection}{Lee Entidad}
\subsection*{Lee Entidad}

\phantomsection
\addcontentsline{toc}{subsection}{Escribe Cabecera Entidades}
\subsection*{Escribe Cabecera Entidades}

\phantomsection
\addcontentsline{toc}{subsection}{Escribe Entidad}
\subsection*{Escribe Entidad}

\phantomsection
\addcontentsline{toc}{subsection}{Inserta Entidad}
\subsection*{Inserta Entidad}

\phantomsection
\addcontentsline{toc}{subsection}{Consulta Entidades}
\subsection*{Consulta Entidades}

\phantomsection
\addcontentsline{toc}{subsection}{Baja Entidad}
\subsection*{Baja Entidad}

\phantomsection
\addcontentsline{toc}{subsection}{Elimina Entidad}
\subsection*{Elimina Entidad}

\phantomsection
\addcontentsline{toc}{subsection}{Reescribe Entidad}
\subsection*{Reescribe Entidad}

\phantomsection
\addcontentsline{toc}{subsection}{Modifica Entidad}
\subsection*{Modifica Entidad}
\newpage

\section{Diagramas}

\newpage
\section{Archivo}

\newpage
\section{Clave primaria}

\newpage
\section{Conclusiones}

\newpage
\phantomsection
\section*{Referencias}
\addcontentsline{toc}{section}{Referencias}

\end{document}